%% kapa: exam=1 round=29 year=2018 period=1 form=A subject=economics subject_ko=경제학원론 seq=5 number=45
%% source: Q-Net 감정평가사 기출문제 (cst00309) artlSeq=5202808 file=extracted/1008/5202808/2018년 제29회 감정평가사 1차시험 기출문제/2018년 제29회 감정평가사 1차시험 기출문제/1교시 A형.pdf
%% conversion: kordoc 4.24.51 -> tools/gen1.py (auto); review=proofread-sonnet; flags=tikz
\begin{question}{45}
기업 A의 총비용곡선에 관한 설명으로 옳지 않은 것은? (단, 생산요소는 한 종류이며, 요소가격은 변하지 않는다.)
\begin{center}
\begin{tikzpicture}[x=0.12mm,y=-0.12mm,font=\small,line width=0.8pt]
\draw[-{Stealth[length=1.6mm]}] (68,305) -- (68,10);
\draw[-{Stealth[length=1.6mm]}] (68,305) -- (362,305);
\draw[line width=0.9pt] plot[smooth] coordinates {(68,248) (80,225) (100,198) (130,175) (170,158) (220,145) (270,130) (310,110) (330,85) (346,45)};
\node[anchor=east] at (62,28) {총비용};
\node at (338,322) {생산량};
\node at (52,322) {0};
\end{tikzpicture}
\end{center}
\begin{choices}
  \item 총평균비용곡선은 U자 모양을 가진다.
  \item 총평균비용이 하락할 때 한계비용이 총평균비용보다 크다.
  \item 평균고정비용곡선은 직각 쌍곡선의 모양을 가진다.
  \item 생산량이 증가함에 따라 한계비용곡선은 평균가변비용곡선의 최저점을 아래에서 위로 통과한다.
  \item 생산량이 증가함에 따라 총비용곡선의 기울기가 급해지는 것은 한계생산이 체감하기 때문이다.
\end{choices}
\end{question}
\answer{2}
